\documentclass[12pt]{article}

% Packages
\usepackage[margin=1in]{geometry}
\setlength{\headheight}{41.55003pt}
\addtolength{\topmargin}{-29.55003pt}
\usepackage{fancyhdr}
\usepackage{datetime}
\usepackage{amsmath}
\usepackage{float}
\usepackage[colorlinks=true, linkcolor=black]{hyperref}

% Date format: MM/DD/YYYY
\newdateformat{mydate}{\twodigit{\THEMONTH}/\twodigit{\THEDAY}/\THEYEAR}

% Header setup
\pagestyle{fancy}
\fancyhf{}
\renewcommand{\headrulewidth}{0pt}
\rhead{Brent Myers \\ Econ 5253 \\ \mydate\today}

\begin{document}

\begin{center}
\textbf{\underline{Problem Set 9}} \\
\vspace{4pt}
\textit{Cross-Validation Tuning of LASSO, Ridge, and Elastic Net Prediction Models}
\end{center}

\section*{Data and Preprocessing}

The data used in this analysis is the Boston Housing dataset from UC Irvine, loaded via the \texttt{mlbench} package in R. The dataset contains 506 observations and 14 variables describing characteristics of housing in the Boston area. The outcome variable is \texttt{medv}, the median value of owner-occupied homes (in \$1,000s), and the remaining 13 variables serve as predictors.

After splitting the data into 80\% training (404 observations) and 20\% testing (102 observations), a preprocessing recipe was applied. The recipe log-transformed the outcome variable, converted the \texttt{chas} dummy variable to a factor, created interaction terms among all continuous predictors, and generated 6th-degree polynomial terms for each continuous predictor. This expanded the feature space from 13 original predictors to 74 predictor (61 more X variables) columns in the training data (75 total columns including the outcome).

\section*{LASSO Regression}

A LASSO model ($\alpha = 1$) was estimated using \texttt{cv.glmnet} with 6-fold cross-validation to tune the penalty parameter $\lambda$. LASSO penalizes the sum of absolute values of the coefficients, which allows it to shrink some coefficients exactly to zero, effectively performing variable selection.

The optimal $\lambda$ selected by cross-validation was 0.002646. The in-sample RMSE was 0.1394 and the out-of-sample RMSE was 0.1831.

\section*{Ridge Regression}

A ridge regression model ($\alpha = 0$) was estimated using the same 6-fold cross-validation procedure. Ridge penalizes the sum of squared coefficients, which shrinks all coefficients toward zero but never eliminates any variable entirely.

The optimal $\lambda$ selected by cross-validation was 0.03664. The in-sample RMSE was 0.1404 and the out-of-sample RMSE was 0.1820.

\section*{Results}

\autoref{tab:results} below shows a comparison between the LASSO and Ridge models. Ridge regression achieved a lower out-of-sample RMSE (0.1820 vs.\ 0.1831), suggesting it generalizes slightly better to unseen data, despite LASSO having a marginally lower in-sample RMSE.

\begin{table}[H]
\centering
\caption{Cross-Validation Results: LASSO and Ridge Regression}
\label{tab:results}
\begin{tabular}{lccc}
\hline\hline
Model & Optimal $\lambda$ & In-Sample RMSE & Out-of-Sample RMSE \\
\hline
LASSO ($\alpha = 1$)  & 0.002646 & 0.1394 & 0.1831 \\
Ridge ($\alpha = 0$)  & 0.03664  & 0.1404 & 0.1820 \\
\hline\hline
\end{tabular}
\end{table}


\section*{Discussion}

\subsection*{Could you estimate OLS with more columns than rows?}

No. OLS requires computing $\hat{\beta} = (X'X)^{-1}X'Y$, which involves inverting the matrix $X'X$. When there are more predictors than observations ($p > n$), this matrix is singular and cannot be inverted. LASSO and ridge regression solve this problem by adding a penalty term to the objective function, which makes estimation possible even in high-dimensional settings. In our case, we have 74 predictors and 404 training observations, so OLS would still be feasible. However, with that many variables relative to the sample size, OLS would be prone to overfitting.

\subsection*{Bias-Variance Tradeoff}

The results illustrate the bias-variance tradeoff. LASSO achieved a slightly lower in-sample RMSE (0.1394) compared to ridge (0.1404), meaning LASSO fit the training data more closely. However, ridge achieved a lower out-of-sample RMSE (0.1820 vs.\ 0.1831), indicating better generalization to unseen data.

This pattern is characteristic of the bias-variance tradeoff: a more flexible model (lower bias) tends to have higher variance, leading to a larger gap between in-sample and out-of-sample performance. Ridge's stronger regularization introduces slightly more bias but reduces variance, resulting in better predictions on the test set. The cross-validation procedure helps find the $\lambda$ that optimally balances this tradeoff for each model.

\section*{Additional Analysis: Elastic Net}

As an additional exercise beyond what was required, an elastic net model ($\alpha = 0.5$) was also estimated. Elastic net combines the LASSO and ridge penalties, blending variable selection with coefficient shrinkage. The problem set introduction references this model, defining $\alpha \in (0,1)$ as the elastic net case.

The optimal $\lambda$ was 0.005807, with an in-sample RMSE of 0.1402 and an out-of-sample RMSE of 0.1803. Interestingly, elastic net achieved the lowest out-of-sample RMSE of all three models, suggesting that a blend of both penalty types may provide the best bias-variance balance for this dataset.

\begin{table}[H]
\centering
\caption{Cross-Validation Results: All Three Models (Including Elastic Net)}
\label{tab:enet_results}
\begin{tabular}{lccc}
\hline\hline
Model & Optimal $\lambda$ & In-Sample RMSE & Out-of-Sample RMSE \\
\hline
LASSO ($\alpha = 1$)        & 0.002646 & 0.1394 & 0.1831 \\
Elastic Net ($\alpha = 0.5$) & 0.005807 & 0.1402 & 0.1803 \\
Ridge ($\alpha = 0$)        & 0.03664  & 0.1404 & 0.1820 \\
\hline\hline
\end{tabular}
\end{table}


These results suggest that for this dataset, all three regularization methods perform similarly, with elastic net providing a marginal improvement.


\end{document}
